\subsection{Ingestion counts}\label{subsec:ingestion-counts}

The main corpus, which the document selection for the document and knowledge extraction stages are made from, is derived from the initial filtering of biomedical papers, described in Subsection~\ref{subsec:corpus-source-and-scope}. The paper counts for the evaluation of different sections should not be conflated.

The papers downloaded from PMC OA based on the filtering, and the counts before and after filtering, are listed in Table~\ref{tab:corpus-acquisition-counts}. 

\begin{table}[htpb]
  \caption[Corpus acquisition and filtering counts]{Corpus acquisition counts. Distinguishes downloaded PDF files, targeted PMC identifiers, file counts after the filters, and successfully ingested papers.}
  \label{tab:corpus-acquisition-counts}
  \centering
  \begin{tabular}{l l r}
    \toprule
      Count & Unit & Value \\
    \midrule
        Downloaded PDF files & File on disk & \textbf{1\,132} \\
        Targeted PMC identifiers & Identifier list & \textbf{1\,093} \\
        Single-PDF papers after supplementary-PDF exclusion & Paper & \textbf{1\,070} \\
        Candidates after \(\leq 30\)-page filter & Paper & \textbf{981} \\
        Successfully ingested papers in the database & Document row & \textbf{977} \\
    \bottomrule
  \end{tabular}
\end{table}

The acquisition step targets the 1\,093 targeted PMC OA identifiers selected in Subsection~\ref{subsec:corpus-source-and-scope}. The 1\,132 downloaded PDF files are not a paper count. This number is higher than the number of targeted PMC OA identifiers, as some PMC packages contain supplementary PDF files in addition to the main articles. 

Two filters define the main corpus selected for document extraction. First, any package that contains supplementary PDF material is excluded entirely from the corpus; only papers with a single, self-contained PDF text are kept, leaving 1\,070 papers. Second, a page count cap for faster document extraction is run, and files with more than 30 pages (p92 when page counts are in ascending order) are filtered out: the long tail is a small subset of papers with a large share of extraction time, so capping the page counts roughly halves extraction time, while dropping under a tenth of papers, as seen in Table~\ref{tab:page-cap-tradeoff}. This leaves 981 candidates, of which 977 are parsed successfully, extracted and ingested into the database; the details of the extraction is explained in Section~\ref{section:Document_Extraction}. Four candidate papers fail Docling conversion on malformed PDF files (missing page-dimension metadata). 

\begin{table}[t]
  \centering
  \caption{Page-count cap trade-off over the single-PDF candidate pool ($n=1\,070$ papers, 16\,277 pages total). Page-volume is a proxy for document-extraction cost.}
  \label{tab:page-cap-tradeoff}
  \begin{tabular}{@{}r r r r r@{}}
    \toprule
    Cap & Papers & \% papers & Page-volume & \% volume \\
    \midrule
    $\le 15$ & 876 & 81.9 & 6\,757 & 41.5 \\
    $\le 20$ & 932 & 87.1 & 7\,731 & 47.5 \\
    $\le 25$ & 963 & 90.0 & 8\,436 & 51.8 \\
    \textbf{$\le 30$} & \textbf{981} & \textbf{91.7} & \textbf{8\,929} & \textbf{54.9} \\
    $\le 40$ & 1\,000 & 93.5 & 9\,602 & 59.0 \\
    $\le 50$ & 1\,014 & 94.8 & 10\,237 & 62.9 \\
    none & 1\,070 & 100.0 & 16\,277 & 100.0 \\
    \bottomrule
  \end{tabular}
\end{table}

The 977 ingested papers therefore are the final corpus for the downstream stages, explained in Sections~\ref{section:Document_Extraction} and~\ref{section:Knowledge_Extraction}. The other counts are reported for transparency: they describe acquisition inputs before and after the two discussed filters. The single-PDF and \(\leq 30\) page selection is reproducible from the extraction CLI; the mechanics of fetching, unpacking, and organizing the files are steps documented in the repository.